% Transcription — fonds Alexandre Grothendieck, shelfmark 29, batch 7 (pages 121-140).
% Established from the facsimile archives/batches/29/batch-07.pdf (a mirror of
% https://grothendieck.umontpellier.fr/29.pdf), per the skill
% transcribe-grothendieck.
%
% Pass: Opus 5 (claude-opus-5), 2026-08-29 — first pass, unchecked against the
% pages by a human.
%
% Fifteen of the twenty pages are transcribed. Five are not:
%
%   - page 123, an empty folder cover inscribed, in his hand, « "Vrai" groupe
%     fondamental [feuilles anciennes] ». It is the wrapper of the run of
%     handwritten sheets that closes the batch, and names it; as a separator
%     sheet it gets no \page{}, exactly as page 111 of batch 6;
%   - pages 131, 133 and 135, three sheets of an English typescript on
%     k-structures on vector spaces (fields of definition, Proposition 1 on
%     k(W), a section "5. Conjugates"), paginated I-6, I-5, I-4 — i.e. in
%     descending order, since the handwritten pages 130, 132 and 134 are
%     written on their backs. They carry nothing in his hand, which is the
%     criterion batches 5 and 6 used for typed versos: kept when annotated,
%     skipped when not. Page 131 is scanned upside down for the same reason.
%
% Page 139 is kept although it carries a single line. That line — "Soit
% (R, λ) le produit fibré de (P, ξ) et (P', ξ') sur (Q, η)" — is the sentence
% page 140 restarts and this time carries through: the page is an abandoned
% attempt, and its abandonment is the record of where the argument was
% re-founded.
%
% The batch falls into three runs.
%
% Pages 121-122 finish the SGA 7 typescript on the functorial behaviour of
% inertia groups that batch 6 broke off inside its § 5. They carry the
% argument to its point: for x̄ varying in a component D_i of the regular
% locus of D of characteristic ≠ ℓ, the inertia homomorphism Z_ℓ(1)(x̄) → G
% is well defined up to conjugation as soon as the Tate sheaf Z_ℓ(1) is
% trivial on D(ℓ), and the classes so obtained are compatible with a morphism
% f : X → Y of the kind studied in the earlier numbers.
%
% Pages 124-129 are a six-page typed letter to Serre, dated 18.X.1959, on the
% "infinitesimal" part of the fundamental group. It is the theoretical heart
% of the batch. Grothendieck sets up an auxiliary category C, a category G of
% C-groups closed under products, kernels of pairs and images, with descending
% chain condition, and an exact functor F from G to group schemes over S; the
% set Z(S, a; G) of pointed principal homogeneous bundles of group F(G) is
% then a functor on G commuting with products and with kernels of pairs, and
% the minimal couples (G, z) form a filtered projective system whose limit
% he writes π₁^G(S, a) — over a base field, π₁(S/k, a), "le groupe
% proalgébrique (mais avec partie infinitésimale) fondamental du k-schéma S".
% He conjectures the exactness relating π₁(X), the inverse image of π₁(Y) and
% the fundamental groups of the fibres of a proper morphism with f_*(O_X) =
% O_Y, computes π₁(X/k) = lim← ₙX for an abelian variety, notes the Cartier
% duality with the ind-algebraic limit of the ₙX*, and closes with the
% timetable of the multiplodocus.
%
% Pages 130 to 140 are the handwritten sheets the folder cover calls "feuilles
% anciennes". They are earlier or parallel workings of the same theory, in his
% own hand and with his own notation: the ℨ of these pages is the Z(S, a; G)
% of the letter. They are hard reading — faint ink, heavy deletion, several
% restarts — and the apparatus here is correspondingly heavy. Page 130 sets up
% the functor F_G(T) of pointed principal bundles on X_T; page 132 lists five
% conditions (i)-(v) on a functor F on (Sch)/S guaranteeing representability by
% an étale separated quasi-finite S-scheme; page 134 states a gluing lemma for
% a compatible family of X(s) over the local rings; pages 136-140 define the
% admissible S-groups by the exact sequence e → U → G → Γ → e and prove, twice
% over, that u ↦ u_*(α) is injective on Hom^gr_S(G, G') when α is minimal.
%
% Notation held for the batch: Z_ℓ(1) underlined as he types it, as in batch 6;
% the categories of the letter kept underlined (C, G) as typed; the gothic ℨ of
% the handwritten pages set as \mathfrak{Z} and flagged once against the
% letter's Z(S, a; G).

\documentclass[11pt,a4paper]{article}
% Le préambule commun à toutes les transcriptions du fonds Grothendieck.
%
% Il tient en une idée : **l'incertitude se compose, elle ne se résout pas.**
% Une transcription qui lisse un mot douteux a détruit la seule chose pour
% laquelle on la faisait — un lecteur doit pouvoir savoir, à chaque endroit,
% ce qui était sur le feuillet et ce qui a été ajouté. Les six macros qui
% suivent sont donc l'appareil critique, et non de la décoration.
%
% Les mêmes commandes sont reconnues par `scripts/render.mjs`, qui produit la
% vue de lecture du site. Une macro ajoutée ici doit l'être aussi là-bas,
% faute de quoi le rendu échouera bruyamment — ce qui est voulu : un rendu
% silencieusement faux est pire qu'un rendu absent.

\ProvidesPackage{grothendieck}[2026/08/08 Transcriptions du fonds Grothendieck]

\usepackage{amsmath,amssymb,amsthm}
\usepackage{mathrsfs}
\usepackage{stmaryrd}
% Les diagrammes commutatifs sont partout dans ce fonds : c'est la langue
% même de la géométrie algébrique de Grothendieck, et une transcription qui
% les rendrait en prose ne transcrirait rien.
\usepackage{tikz-cd}
% Français par défaut — c'est la langue des feuillets — mais l'anglais est
% chargé aussi : la traduction et le résumé sont des documents anglais, et la
% typographie française (espaces avant les deux-points) y serait une faute.
% Ces documents basculent par \selectlanguage{english} en tête de corps.
\usepackage[english,french]{babel}
\usepackage{fontspec}
\usepackage{geometry}
\geometry{a4paper,margin=25mm,marginparwidth=28mm}
\usepackage{marginnote}
\usepackage{xcolor}
% ulem, not soul. `soul`'s \st and \ul are built on a character-by-character
% scan that XeLaTeX's font machinery breaks: the document fails with an
% undefined control sequence rather than degrading. ulem does the same two jobs
% and is Unicode-engine safe. `normalem` stops it hijacking \emph, which the
% transcriptions use for Grothendieck's own emphasis.
\usepackage[normalem]{ulem}
\usepackage{hyperref}
\hypersetup{colorlinks=true,linkcolor=black,urlcolor=black,pdfborder={0 0 0}}

% --- Métadonnées du lot -----------------------------------------------------
% Reprises telles quelles par le script de rendu, qui les affiche en tête de
% la vue de lecture. Elles fixent ce que le fichier prétend couvrir : sans
% elles, une transcription flotte sans point d'attache dans un fonds de seize
% mille feuillets.

\newcommand{\folder}[1]{\gdef\grothFolder{#1}}
\newcommand{\batch}[1]{\gdef\grothBatch{#1}}
\newcommand{\leaves}[2]{\gdef\grothFirst{#1}\gdef\grothLast{#2}}
\newcommand{\foldertitle}[1]{\gdef\grothFoldertitle{#1}}
\gdef\grothFolder{?}\gdef\grothBatch{?}\gdef\grothFirst{?}\gdef\grothLast{?}\gdef\grothFoldertitle{}

% --- Le résumé --------------------------------------------------------------
%
% Il ouvre la lecture modernisée, sous le titre « Résumé », et il est composé
% en petit. Ce n'est pas un ornement : c'est le seul passage du document qui
% ne soit pas des mathématiques, et un lecteur doit pouvoir le reconnaître
% avant de l'avoir lu — puis le sauter s'il connaît déjà le sujet, ou s'y
% arrêter s'il ne le connaît pas. Le filet qui le suit marque la reprise du
% corps.

\newenvironment{resume}
  {\small\setlength{\parskip}{0.45em}}
  {\par\medskip\hrule\medskip}

% --- Mots-clés ---------------------------------------------------------------
%
% En anglais, en clôture du résumé de la lecture modernisée : le vocabulaire
% moderne sous lequel chercher ce que les feuillets construisent. Le manifeste
% du site (`npm run manifest`) les extrait pour étiqueter la cote entière —
% cette ligne est la source unique des tags, il n'y a pas de fichier à côté.
\newcommand{\keywords}[1]{\par\smallskip\noindent
  {\footnotesize\textsc{Keywords} --- \textit{#1}}\par}

% --- L'appareil critique ----------------------------------------------------

% Le numéro de feuillet, en marge. C'est l'ancre qui permet de régler un
% désaccord sur une formule en regardant *un* feuillet plutôt qu'une chemise
% de six cents. Le site s'en sert aussi pour tourner les pages du fac-similé
% quand on fait défiler la transcription.
\newcommand{\leaf}[1]{%
  \marginnote{\footnotesize\color{gray!70}\textsf{#1}}%
  \label{leaf:#1}\ignorespaces}

% Illisible : on ne devine pas. Un mot inventé qui se lit comme les autres est
% le pire résultat possible.
\newcommand{\ill}{\textcolor{red!60!black}{[\,\ldots\,]}}

% Lecture douteuse : le mot est donné, et signalé comme incertain.
\newcommand{\uncertain}[1]{\uline{#1}}

% Ajout de l'éditeur — une lettre manquante, un mot sous-entendu.
\newcommand{\add}[1]{\textcolor{blue!50!black}{[#1]}}

% Biffé par Grothendieck lui-même. Ce qu'il a rayé fait partie du document :
% une rature dit souvent où la pensée a changé de direction.
\newcommand{\struck}[1]{\sout{#1}}

% Note du transcripteur, sur l'état matériel du feuillet.
\newcommand{\note}[1]{\par\smallskip{\small\color{gray!60!black}\textit{[#1]}}\par\smallskip}

% Note marginale de Grothendieck, distincte du corps du texte.
\newcommand{\marginal}[1]{\par\smallskip\noindent\hspace{1em}%
  {\small\color{gray!60!black}#1}\par\smallskip}

% --- Titre ------------------------------------------------------------------

\AtBeginDocument{%
  \begin{center}
    {\large\bfseries Fonds Alexandre Grothendieck --- cote n\textsuperscript{o}~\grothFolder}\\[2pt]
    {\itshape\grothFoldertitle}\\[4pt]
    {\small feuillets \grothFirst--\grothLast\ (lot \grothBatch)}\\[2pt]
    {\footnotesize\color{gray!60!black}Transcription établie d'après le fac-similé numérisé
     par l'Université de Montpellier.\\
     Les lectures incertaines sont soulignées, les illisibles notées [\,\ldots\,],
     les ajouts entre crochets.}
  \end{center}
  \vspace{1em}}

\endinput


\folder{29}
\batch{7}
\pages{121}{140}
\dating{1959-1969}
\watermark{Édition de démonstration}
\foldertitle{Groupe fondamental [Autour de SGA 4 et SGA 7] : lettres (1959, 1967, 1969), tapuscrits et copies de tapuscrit annotés (s.d.), notes manuscrites (s.d.)}

\begin{document}

\section*{Groupes d'inertie : classes de conjugaison d'homomorphismes de $\underline{Z}_{\ell}(1)$ (pages 121 à 122)}

\page{121}

$X$ (donc \struck{et} $X$ est régulier en les points de son support),
\struck{$U = f^{-1}(V) = X - \operatorname{supp} D$}
\marginal{(Condition de transversalité)}

\begin{enumerate}
\item[b)] $\operatorname{supp} E$ est irréductible, et il est régulier en les
points de $f(\operatorname{supp} D) = \operatorname{supp} E \cap f(X)$ (donc $Y$
est régulier en ses points).
\end{enumerate}

Alors pour tout $x \in \operatorname{supp} D$, \struck{les images des} les
sous-groupes d'inertie de $G$ en $x$ sont comme images par $\varphi : G \to H$
des sous-groupes de $H$ qui sont des sous-groupes d'inertie de $H$
relativement au point maximal de $\operatorname{supp} E$. Par suite, pour toute
représentation linéaire de $G$ (dans un module $M$ sur quelque anneau)
\struck{\ill{}} qui veut bien se factoriser par une représentation linéaire de
$H$, les images dans $\operatorname{Aut}(M)$ des sous-groupes d'inertie
\struck{dans} de $G$ relatifs aux divers points de $\operatorname{supp} D$ (ou
encore, relatifs aux diverses composantes irréductibles de $\operatorname{supp}
D$, cela revient au même en vertu de 3) sont conjugués dans
$\operatorname{Aut}(M)$. \marginal{Remords au n\textsuperscript{o} 3, 4.}

\struck{Application} On peut demander \struck{du} un résultat plus précis, en
considérant pour tout $\xi$ comme ci-dessus (au dessus d'un point géométrique
$\bar{x}$ de $\operatorname{supp} D$) l'homomorphisme correspondant

\[ \underline{Z}_{\ell}(1)(\bar{x}) \to G \]

dont l'image est le sous-groupe d'inertie correspondant, et en se demandant
comment cet homomorphisme varie avec le choix de $\xi$ et de $\bar{x}$. Pour
$\xi$ variable, la variation se fait évidemment modulo composition avec les
automorphismes intérieurs de $G$, \struck{Pour donner un sens à la question}
donc pour $\bar{x}$ fixé, on trouve une classe de conjugaison d'homomorphismes

\[ \underline{Z}_{\ell}(1)(\bar{x}) \to G. \]

\struck{pour} \add{Étudions} la variation de cet homomorphisme, pour $\bar{x}$
variable dans une composante \struck{de la} \add{$D_{i}$ de} l'ensemble des
points réguliers de $D$ \add{à car.\ $\neq \ell$} ; d'après ce qu'on a vu au
n\textsuperscript{o} 3, on passe d'un tel homomorphisme à un autre en composant
avec un quelconque des isomorphismes $\underline{Z}_{\ell}(1)(\bar{x}) \to
\underline{Z}_{\ell}(1)(\bar{x}')$ associé à une classe de chemins de $\bar{x}$
à $\bar{x}'$. Si cet \uncertain{isomorphe}\note{sic ; lire
« isomorphisme »} ne dépend pas du choix de cette classe de chemins, i.e.\ si
\struck{l'action du groupe fondamental} le faisceau $\ell$-adique de Tate
$\underline{Z}_{\ell}(1)$ sur $D(\ell)$ est trivial (ce qui se voit déjà sur la

\page{122}

restriction dudit faisceau au \struck{corps résiduel} point générique de
$D_{i}$), alors on peut identifier entre eux tous les groupes
$\underline{Z}_{\ell}(1)(\bar{x})$ envisagés à un même groupe
$\underline{Z}_{\ell}(1)(D_{i})$ (qu'on peut interpréter comme groupe des
sections de $\underline{Z}_{\ell}(1)$ sur $D_{i}$ ou sur le point générique de
$D_{i}$), et on trouve une classe \add{de conjugaison} \struck{bien déterminée}
d'homomorphismes

\[ \underline{Z}_{\ell}(1)(D_{i}) \to G. \]

(On trouve un résultat analogue pour les homomorphismes de produits de groupes
$\underline{Z}_{\ell}(1)$ envisagés dans le n\textsuperscript{o} 3.) Bien
entendu, si $\underline{Z}_{\ell}(1)$ est trivial sur l'ouvert \add{$X_{\ell}$}
des points de $X$ \struck{lorsque l'on a un $f : X \to Y$ comme aux
n\textsuperscript{os} \ill{}} de car.\ $\neq \ell$, et si ce dernier est
connexe, on peut identifier les groupes $\underline{Z}_{\ell}(1)(D_{i})$ au
même groupe $\underline{Z}_{\ell}(1)(X_{\ell})$, \struck{si on} soit
$\underline{Z}_{\ell}(1)$, et on trouve \struck{un homomorphisme dans} pour
chaque $i$ une \struck{classe} classe de conjugaison d'homomorphismes
$\underline{Z}_{\ell}(1) \to G$. En général, celles-ci \struck{sont} peuvent
être distinctes pour des $i$ distincts, bien entendu. Mais relativement à un
morphisme $f : X \to Y$ comme dans les n\textsuperscript{os} 2, 4, la formation
des classes \add{de conjugaison} \struck{précédentes} d'homomorphismes
\add{ci-dessus} a un caractère fonctoriel évident, que le lecteur explicitera
lui-même s'il en a besoin.

\section*{Lettre à Serre : la partie « infinitésimale » du groupe fondamental (pages 124 à 129)}

\page{124}

\uncertain{Vallocin}\note{le nom de lieu est lu sans certitude}, 18.X.1959

Mon cher Serre,

Tate m'a écrit de son côté \struck{\ill{}} sur ses histoires de courbes
elliptiques, et pour me demander si j'avais des idées sur une définition
globale des variétés analytiques sur des corps valués complets. Je dois avouer
que je n'ai pas du tout compris pourquoi ses résultats suggèreraient
l'existence d'une telle définition, et suis encore sceptique. Je n'ai pas non
plus l'impression d'avoir rien compris à son théorème, qui \struck{\ill{}} ne
fait que exhiber par des formules brutales un certain isomorphisme de groupes
analytiques ; \struck{\ill{}} on conçoit que d'autres formules tout aussi
explicites en donneraient un autre pas plus mauvais (sauf preuve du contraire !)

J'ai réfléchi un petit peu à la partie ``infinitésimale'' du groupe
fondamental, juste assez pour me convaincre que ça existe et est raisonnable.
Voici un contexte (certainement insuffisant d'ailleurs) où ça marche. On a un
schéma $S$ \add{(connexe)} (par exemple un schéma algébrique sur un corps $k$),
une catégorie auxiliaire $\underline{C}$ (par exemple la catégorie des schémas
algébriques finis sur $k$), une catégorie \struck{\ill{}} $\underline{G}$ dont
les objets sont des $\underline{C}$-groupes, et les morphismes des morphismes
de $\underline{C}$-groupes (par exemple les groupes algébriques finis sur $k$).
On suppose que dans $\underline{C}$ les produits fibrés existent, et que
$\underline{G}$ satisfait aux conditions suivantes :

\begin{enumerate}
\item[(i)] $\underline{G}$ est stable par produits, et si $u, v : G \to G'$
sont des morphismes dans $\underline{G}$, alors le noyau du couple $(u,v)$
i.e.\ le sous-groupe \struck{\ill{}} maximum de $G$ où $u$, $v$ coïncident ---
qui existe, car s'exprime par un produit fibré --- est dans $\underline{G}$.
\item[(ii)] Si $u : G \to G'$ est un morphisme dans $\underline{G}$,
\struck{\ill{}} le groupe image existe, est isomorphe à un quotient de $G$
comme il se doit, et est dans $\underline{G}$.
\item[(iii)] Toute suite décroissante de sous-groupes \struck{\ill{}} $\in
\underline{G}$ d'un $G \in \underline{G}$ est stationnaire.
\end{enumerate}

On suppose donné de plus un foncteur covariant \struck{\ill{}} $F$ de
$\underline{G}$ dans les \struck{groupes} schémas en groupes sur $S$ (par
exemple $G \to S \times_{k} G$), et on suppose :

\begin{enumerate}
\item[(iv)] Le foncteur $F$ commute aux produits, \add{aux} noyaux de couples
de morphismes, images (on pourra dire que $F$ est ``exact'').
\item[(v)] Si $H$ est \struck{\ill{}} de la forme $F(G)$ ($G \in
\underline{G}$) alors $H$ \add{est ``spécial'' i.e.} admet une suite exacte
\end{enumerate}

\page{125}

de schémas en groupes finis et plats sur $S$

\[ e \to H_{\mathrm{inf}} \to H \to H_{\mathrm{sép}} \to e \]

où $H_{\mathrm{inf}}$ est purement infinitésimal (i.e.\ la projection
$H_{\mathrm{inf}} \to S$ est géométriquement injectif) et où
$H_{\mathrm{sép}}$ est non ramifié sur $S$. (Dans le cas d'un corps de base
$k$, une telle suite exacte est déduite d'une suite exacte analogue pour un
groupe algébrique fini $G$ sur $k$ ; d'ailleurs si $k$ est parfait, cette suite
splitte canoniquement, car $G_{\mathrm{réd}}$ est \add{alors} un sous-groupe de
$G$, isomorphe à $G_{\mathrm{sép}}$ par la projection $G \to
G_{\mathrm{sép}}$. Faire attention cependant que $G_{\mathrm{réd}}$ va opérer
sur $G_{\mathrm{inf}}$, on aura seulement un produit semi-direct.) Enfin
\struck{\ill{}}

\begin{enumerate}
\item[(vi)] $S$ est réduit.
\end{enumerate}

Les conditions (v) et (vi) ont une sale gueule, et sont essentiellement
provisoires. Elles servent à \struck{assurer la validité} \add{causer} du

\textbf{Lemme.} Soit $H$ comme dans (v), $S$ comme dans (vi), $P$ un fibré
principal homogène sous $H$, $Q$ un autre, $u$, $v$ deux isomorphismes de $P$
dans $Q$ transformant un ``point \struck{\ill{}} \add{marqué}'' donné en un
autre donné, alors $u = v$.

(On est ramené, en tordant $H$, au cas où $P$ est trivial, et alors cela
résulte du fait suivant : une section de $Q$ est un isomorphisme de $S$ sur une
composante connexe de $Q_{\mathrm{réd}}$.)

Remarque néanmoins que les conditions (i) à (vi) n'excluent pas des groupes
structuraux tordus (relativement à un corps de base $k$).

Soit maintenant $a$ un ``point marqué'' de $S$ (i.e.\ une extension
algébriquement close du corps résiduel d'un $s \in S$). Pour tout $G \in
\underline{G}$, on désigne par $Z(S, a ; G)$ ou simplement $Z(G)$ l'ensemble
des classes (à isomorphisme près) de fibrés principaux homogènes sur $S$, de
groupe $F(G)$, munis d'un point marqué au dessus de $a$. Évidemment $Z(G)$ est
un foncteur de $\underline{G}$ dans la catégorie des ensembles. Ce foncteur
satisfait aux conditions suivantes :

\begin{enumerate}
\item[1)] Il commute aux produits (\struck{\ill{}} car $F$ y commute).
\item[2)] Il commute aux noyaux de couples, en d'autres termes si $G'' \to G
\rightrightarrows G'$ est exact dans $\underline{G}$, alors $Z(G'') \to Z(G)
\rightrightarrows Z(G')$ est exacte, i.e.\ $Z(G'')$ s'identifie à l'ensemble
des éléments de $Z(G)$ dont l'image dans $Z(G')$ par $Z(u)$ et $Z(v)$ est la
même. (Ceci résulte de l'exactitude de $F$, et du lemme.)
\end{enumerate}

\page{126}

De ces deux propriétés résulte formellement qu'on peut trouver un système
projectif \add{filtrant} $(G_{i})_{i \in I}$ dans $\underline{G}$, avec des
morphismes $G_{i} \to G_{j}$ qui sont des épimorphismes dans $\underline{C}$,
``essentiellement unique'' dans un sens facile à préciser, tel que
\struck{\ill{}} l'on ait un isomorphisme fonctoriel

\[ Z(G) = \varinjlim \operatorname{Hom}_{\underline{G}}(G_{i}, G) \]

C'est ce système projectif (pris modulo une équivalence qui intuitivement
signifie qu'on a passé à la limite projective des $G_{i}$) qui peut être noté
$\pi_{1}^{\underline{G}}(S, a)$ et joue le rôle d'un groupe fondamental de $S$
en $a$ (relativement à la catégorie $\underline{G}$ de groupes, munie du
foncteur $F$). Dans le cas \struck{\ill{}} d'un corps de base $k$,
$\underline{G}$ étant la catégorie des groupes algébriques finis sur $k$, on
pourra écrire $\pi_{1}(S/k, a)$, c'est le groupe proalgébrique (mais avec
partie infinitésimale) fondamental du $k$-schéma $S$. \marginal{à tout}
Lorsque $k$ est parfait, la décomposition d'un groupe algébrique fini en partie
infinitésimale et partie réduite montre que le groupe fondamental est produit
semi-direct de sa partie réduite ou séparable \struck{\ill{}} (qui correspond à
un groupe compact discontinu ordinaire, sur lequel le groupe fondamental
ordinaire de $k$ --- i.e.\ le groupe de Galois de $\bar{k}$ sur $k$ --- opère),
par sa partie infinitésimale, \struck{\ill{}} qui elle ne dépend d'ailleurs
plus du choix du point base $a$ de $S$. \marginal{(lorsque $a$ \ill{}
automatiquement une \ill{})} Noter que la partie séparable du groupe
fondamental peut se réconstituer facilement à l'aide du groupe fondamental
ordinaire de $S$ et de son homomorphisme naturel dans
$\operatorname{Gal}(\bar{k}/k)$. (\add{mais} Il est ``plus grand'' que le
groupe fondamental ordinaire, parceque il correspond à la classification de
revêtements principaux de groupe structural non seulement un groupe fini
ordinaire, mais un groupe \struck{\ill{}} fini et séparable sur $k$ (i.e.\ un
groupe fini ordinaire sur lequel on fait opérer
$\operatorname{Gal}(\bar{k}/k)$ de façon non nécessairement triviale).)
\struck{\ill{}} J'avoue que si $k$ n'est pas algébriquement clos, le groupe
fondamental ci-dessus n'est \struck{peut-être} \add{sans doute} pas le bon ; on
devrait sans doute prendre le groupe fondamental de $S \times_{k} \bar{k}$, qui
est un groupe proalgébrique défini sur $\bar{k}$, noter qu'il est muni d'une
donnée de descente de $\bar{k}$ à $k$ et est par suite défini en réalité sur
$k$ ; \struck{avec cette définition, si $S = \operatorname{Spec}(k)$, la
\ill{} partie séparable du groupe fondamental de $S$ (relativement à la
ponctuation de $S$ définie par une \ill{} clôture algébrique $\bar{k}$ de
$k$) n'est autre que}

\page{127}

\struck{$\operatorname{Gal}(\bar{k}/k)$, mais où on fait opérer
$\operatorname{Gal}(\bar{k}/k)$ par automorphismes intérieurs. Ainsi, \ill{}
le groupe des points rationnels sur $k$ du groupe fondamental de
$\operatorname{Spec}(k)$ s'identifie au centre de
$\operatorname{Gal}(\bar{k}/k)$.}

Voici comment on voit l'existence du système projectif $(G_{i})$. Un couple
$(G, z)$, avec $G \in \underline{G}$ et $z \in Z(G)$, est dit \emph{minimal} si
on ne peut trouver un sous-groupe $G' \subset G$ de $G$, \struck{\ill{}} $\neq
G$, tel que $z$ soit \struck{dans} de la forme $Z(i)(z')$, avec $z' \in Z(G')$
et $i : G' \to G$ l'injection. \struck{\ill{}} On dit qu'un couple $(G, z)$ est
\emph{majoré} par un couple $(G', z')$ si on peut trouver un homomorphisme
$G' \xrightarrow{\;u\;} G$ tel que $z = Z(u)(z')$. Il résulte de (iii) que tout
couple $(G, z)$ est majoré par un couple minimal, et de la propriété 2) de $Z$
que si $(G', z')$ est minimal et domine $(G, z)$, alors il existe
\emph{un seul} $u : G' \to G$ tel que $z = Z(u)(z')$. De ceci résulte que les
couples $(G, z)$ minimaux forment un système projectif pour la relation de
domination, d'ailleurs filtrant (car $(G, z)$ et $(G', z')$ sont dominés par
$(G \times G', (z, z'))$ lui-même dominé par un $(G'', z'')$ minimal), c'est le
système cherché. On peut si on veut choisir un couple $(G, z)$ dans tout
système de couples \add{minimaux} isomorphes, de \struck{\ill{}} telle façon
que l'ensemble d'indices \struck{\ill{}} $I$ devient ordonné et non seulement
préordonné. (N.B. Ce genre de raisonnements formels est aussi celui qui sert en
théorie des \struck{\ill{}} modules \ldots)

Je ne sais encore comment devrait être formulée la suite exacte d'homotopie,
pour bien faire l'inclusion dans le groupe fondamental d'une partie
infinitésimale devrait donner une théorie satisfaisante du comportement du
groupe fondamental par spécialisation. J'espère que si $f : X \to Y$ est un
morphisme propre et séparable à fibres absolument connexes i.e.\ tel que
$f_{*}(\mathcal{O}_{X}) = \mathcal{O}_{Y}$, $X$ étant muni \struck{\ill{}}
d'une section sur $Y$ qui détermine des points-base sur les fibres
\marginal{(pour simplifier)}, alors les groupes fondamentaux \emph{totaux} des
fibres de $X$ forment sur $Y$ un pro-schéma en groupes \add{filtrant}, de façon
précise qu'on peut trouver un système projectif \add{$\pi_{1}(X/Y, s)$}
$(G_{i})_{i \in I}$ de schémas en groupes $G_{i}$ `spéciaux' sur $Y$, avec des
homomorphismes $G_{i} \to G_{j}$ qui soient des épimorphismes de $Y$-schémas
(i.e.\ correspondant à un homomorphisme injectif de faisceaux cohérents sur
$Y$), de telle façon que les groupes fondamentaux totaux des fibres de $X$
\struck{\ill{}} déduisent dudit pro-schéma en groupes par simple
spécialisation. (C'est en tous cas ce que semble indiquer le cas où $X$ est un
schéma abélien sur $Y$, cf.\ plus bas.)

\page{128}

Lorsqu'il n'y a pas de corps de base, il n'est pas clair cependant comment et à
quels groupes fondamentaux \struck{\ill{}} pour $X$, $Y$ on peut rattacher
\struck{\ill{}} les groupes fondamentaux des fibres, pour tenir lieu de suite
exacte d'homotopie. S'il y a un corps de base, une \add{conjecture}
\struck{\ill{}} de départ serait que \struck{l'image inverse par $f$ de}
$\pi_{1}(X/Y)$ \add{$\pi_{1}(X)$}, le pro-schéma en groupes sur $X$ qui
définit le système local des groupes fondamentaux \struck{\ill{}} de $X$ en ses
divers points, et \struck{enfin} l'image inverse par $f$ du pro-schéma en
groupes analogue \add{$\pi_{1}(Y)$} sur $Y$, sont reliés par une suite exacte
\struck{(qui serait définie d'ailleurs indépendamment de l'existence d'une
section de $X$ sur $Y$, car le premier \ill{} pro-schéma en groupes nommé (sur
$X$) serait défini en tous cas). Si $Y$ est le spectre d'un anneau local
complet, notamment un anneau de valuation discrète, on devrait pouvoir s'y
mettre.} Il y a de la joie en perspective pour les groupes
d'homotopie supérieurs \ldots

Il n'est pas dit que la conjecture envisagée soit plus difficile à démontrer
que la partie déjà connue pour les groupes fondamentaux ordinaires. Elle
impliquerait que pour des schémas complets $X$, $Y$ au dessus d'un corps alg
clos $k$, on a $\pi_{1}(X \times Y/k) = \pi_{1}(X/k) \times \pi_{1}(Y/k)$.
Utilisant tes raisonnements, on en conclut que si $X$ est une variété abélienne
sur $k$, alors $\pi_{1}(X/k)$ est abélien, et qu'un revêtement principal
minimal $X'$ de $X$ est une variété abélienne isogène à $X$. On en déduit

\[ \pi_{1}(X/k) = \varprojlim {}_{n}X \]

où ${}_{n}X$ est le noyau (avec sa partie infinitésimale aussi, bien sûr) de la
multiplication par $n$, l'homomorphisme de ${}_{mn}X$ dans ${}_{n}X$ étant
induit par la multiplication par $m$. (Utilisant Cartier, ce groupe fondamental
est dual, au sens de Cartier, au groupe ind-algébrique limite inductive des
${}_{n}X^{*}$, où $X^{*}$ est la variété duale de $X$.)

Prenant la $p$-composante de ce groupe fondamental, on trouve ce qui, pour le
nombre premier $p$, doit jouer le rôle du module de Weil. Il est hors de doute
(et Cartier doit l'avoir fait) qu'il est possible d'associer \struck{\ill{}}
fonctoriellement à un groupe proalgébrique infinitésimal abélien un module sur
l'anneau de Witt, et \struck{\ill{}} qu'il \add{soit} facile de vérifier qu'en
l'occurence ce module se décompose en les trois morceaux que tu connais bien
(correspondants aux trois types principaux de groupes algébriques
\add{abéliens}), \struck{\ill{}} On trouve ainsi

\page{129}

pour ton théorème une formulation plus naturelle (qui devrait \struck{\ill{}}
fournir une démonstration uniforme, sans distinction du cas $\ell \neq p$ et
$\ell = p$), de même qu'on voit en même temps que ta somme abracadabrante se
comporte bien quand on fait varier $X$ dans une famille, i.e.\ si $X$ est
\struck{\ill{}} un schéma abélien sur $Y$ (car alors les ${}_{n}X$ seront
d'excellents schémas en groupes finis et plats sur $Y$, dont on peut prendre
formellement la limite projective \ldots).

J'espère arriver dans l'année prochaine à une théorie satisfaisante du groupe
fondamental, \struck{\ill{}} et achever la rédaction des Chapitres IV, V, VI,
VII (ce dernier étant le groupe fondamental), en même temps que des catégories.
Dans deux ans résidus, dualité, intersections, Chern, Riemann-Roch. Dans trois
ans cohomologie de Weil, et un peu d'homotopie si Dieu veut. Et entre-temps, je
ne sais quand, le grand théorème d'existence avec Picard etc, un peu de courbes
algébriques, \add{les schémas abéliens.} Sauf difficultés imprévues ou
enlisement, le multiplodoque devrait être fini d'ici 3 ans, ou 4 au maximum. On
pourra commencer à faire de la géométrie algébrique !

Bien à toi,

\section*{Le foncteur des fibrés principaux pointés au-dessus de $X/S$ (page 130)}

\page{130}

$X$ \ill{} muni d'une section $\uncertain{\sigma_{1}}$ ; $f$ propre plat,
$f_{*}(\mathcal{O}_{X}) = \mathcal{O}_{S}$ \emph{universellement}.

\begin{tikzcd}
X \arrow[d, "f"] \\
S
\end{tikzcd}

$G$ désignera un groupe \struck{\ill{}} fini et \struck{séparable}
\uncertain{étale} sur $S$ [le cas important \ill{} celui \ill{} $G = S \times
\Gamma$].

$F_{G}(T) = $ ens.\ des classes d'isomorphie \ill{} de fibrés principaux
homogènes sur $X_{T}$, de groupe $G_{T} \times_{T} X_{T} = (G \times_{S}
X_{T})$, munis d'une section \ill{} le long de $\sigma_{T}$, i.e.\
$\sigma_{T}^{*}(X') \simeq G_{T}$.

N.B. \ill{} automorphismes d'un tel fibré \ill{} \marginal{$= $ pour $G$
commutatif, les \ill{}}

Le foncteur $F_{G}(T)$ en $T$ \ill{} est-il représentable ? Si oui, le \ill{}
qui le représente est étale sur $S$, et \ill{} \marginal{(11.4) \quad 6.4}

Notons que le foncteur $F_{G}$ \ill{} est compatible \ill{} et à la descente
pour les morphismes fid.\ plats \ill{} et quasi-compacts.

\struck{\ill{}} \quad $G$ \quad $H^{1}(X, G)$ \quad $G$ \quad $G$
\uncertain{tordu}\note{fragment isolé, sans lien explicite avec ce qui
précède}

\section*{Un critère de représentabilité par un schéma en groupes étale (page 132)}

\page{132}

$S$ préschéma loc.\ noethérien. $F$ foncteur $(\mathrm{Sch})_{/S} \to
(\mathrm{Ens})$.

\begin{enumerate}
\item[(i)] $F$ compatible aux limites inductives d'anneaux.
\item[(ii)] $F$ compatible aux descentes fidèlement plates quasi-compactes
\add{et localisation}.
\item[(iii)] Si $A$ anneau local noethérien \emph{complet} sur $S$, corps
résiduel $k$, alors $F(A) \simeq F(k)$.
\item[(iv)] $F$ séparé.
\item[(v)] Pour tout corps $k$ sur $S$, $F(k)$ est fini, et si $k'$ est une
ext.\ alg.\ sép.\ de $k$, \struck{\ill{}} alg.\ clos, $F(k) \to F(k')$ est
bijectif.
\end{enumerate}

Alors $F$ représentable par un \struck{\ill{}} \ill{} étale \struck{et gpe}
séparé quasi-fini sur $S$.

\struck{Démonstration.}

Supposons $S$ le spectre d'un anneau local complet $A$, à corps résiduel alg.\
clos, \struck{\ill{}} de dim.\ \ill{}, et \ill{} sur un $S_{1} = S - y$ le
foncteur $F$ soit représentable par \uncertain{$M_{1}$}. \struck{\ill{}}

$\xi_{i}$ \struck{\ill{}} $(i \in I)$ \struck{$\in k$}, éléments de $F(k)$, ils
proviennent d'éléments $\xi_{i} \in F(A)$ bien déterminés, en vertu de (iii) ;
leur restriction $\xi'_{i} = \xi_{i} \mid S_{1} \in F(S_{1})$, ils
correspondent à des \ill{}\note{la page s'arrête ici, la suite manque}

\section*{Recollement des $X(s)$ : lemme d'existence (page 134)}

\page{134}

\textbf{Lemme.} Soit $S$ un préschéma loc.\ noethérien, et pour tout $s \in S$,
soit \struck{\ill{}} $X(s)$ un préschéma quasi-fini étale séparé sur
$\operatorname{Spec}(\mathcal{O}_{s})$, et pour $t \in S$, $s \in \bar{t}$,
soit $\varphi_{st} : X(t) \to X(s)$ rendant commutatif le diagramme ci-contre.

\begin{tikzcd}
X(s) \arrow[d] & X(t) \arrow[l, "\varphi_{st}"'] \arrow[d] \\
\operatorname{Spec}(\mathcal{O}_{s}) & \operatorname{Spec}(\mathcal{O}_{t}) \arrow[l]
\end{tikzcd}

(les $\varphi_{st}$ \ill{} de transitivité) $\varphi_{st}\varphi_{tu} =
\varphi_{su}$. Alors il existe un préschéma \add{quasi-fini séparé} étale sur
$S$, soit $X$, et des isomorphismes

\[ \varphi_{s} : X \times_{S} \operatorname{Spec}(\mathcal{O}_{s})
\xrightarrow{\;\sim\;} X(s), \]

rendant commutatifs \add{des} diagrammes \ill{}. $X$ est unique à isomorphisme
unique près.

Ce dernier fait résulte d'un énoncé plus général sur la catégorie des
préschémas de type fini sur $S$, qui donne un foncteur pleinement fidèle. Je ne
suis pas \ill{} de mon raisonnement \ill{} $S$ \uncertain{affine}, ou même
\uncertain{noethérien}. \ill{} pour \ill{} on regarde l'énoncé \ill{} pour les
sous-préschémas fermés de $S$ $\neq S$.

\marginal{\ill{} pour tout $s \in S$, il existe une \uncertain{réalisation}
\ill{} $X(s)$ sur un voisinage ouvert $U$ de $s$ \ill{} :}

\begin{tikzcd}[column sep=small, row sep=small, nodes={font=\scriptsize}]
X(s) \arrow[d] & X \arrow[l] \arrow[d] & X \times_{U} \operatorname{Spec}(\mathcal{O}_{t}) \arrow[l] \arrow[d] \arrow[r, dashed] & X(t) \\
\operatorname{Spec}(\mathcal{O}_{s}) \arrow[r] & U & \operatorname{Spec}(\mathcal{O}_{t}) \arrow[l] &
\end{tikzcd}

\section*{Groupes admissibles et couples minimaux (pages 136 à 140)}

\page{136}

Soit $S$ un schéma \emph{réduit connexe}.

On considère des schémas en groupes, finis et plats sur $S$, soit $G$, pour
lesquels \ill{} une suite exacte

\[ e \to U \to G \to \Gamma \to e \]

i.e.\ la projection $U \to S$ est plat, fini, \struck{et} et géométriquement
\struck{\ill{} homomorphisme} injectif, et si $\Gamma \to S$ est plat, fini,
séparable. Alors le sous-$S$-groupe $U$ de $G$ est déterminé de façon unique
[pour cette question, l'hypothèse « $S$ réduit » est-elle nécessaire ?].
\marginal{Groupes admissibles ($S$ réduit ?)}

On \struck{va} \add{veut} étudier des fibrés principaux $P$ sous de tels
$S$-groupes. On suppose $S$ pointé par $a \in \uncertain{S}$, et on considère
les $P$ munis d'une structure pointée.

\textbf{Proposition.} \struck{\ill{}} Soit $S' \to S$ un morphisme
\struck{\ill{}}, $G$ un $S$-groupe fini. \add{Si $G \times_{S} S'$ est
admissible, $G$ \uncertain{aussi} ;} \struck{\ill{}} \add{la réciproque est
vraie si $S' \to S$ est fidèlement plat.}

\uncertain{Direct est trivial.} Réciproque résulte du \ill{} théorème de
\ill{} et unicité.

\page{137}

\textbf{Corollaire.} Supposons $G$ admissible, et $P$ un fibré principal sous
$G$. Alors $\operatorname{Aut}(P)$, considéré comme $S$-groupe, est admissible.

\textbf{Proposition 2.} Soit $(P, \xi)$ un fibré principal \emph{pointé} sous
un $S$-groupe admissible $G$. Alors $\operatorname{Aut}(P, \xi)$ est réduit à
l'identité.

Ça résulte, moyennant le corollaire précédent, de \ill{}

\struck{Lemme} \textbf{Proposition 3.} Deux sections d'un \struck{\ill{}}
\add{préschéma} $P$ (\struck{admissible sur $S$}) \add{réduit connexe} qui
\ill{} en un pt, sont identiques. [De façon plus précise, si $s$ est une
section de $P/S$, il existe un sous-préschéma fermé $Q$ de $P$, tel que
$Q_{\mathrm{réd}} \to S$ soit un isom., et l'isomorphisme réciproque induit par
$s$.]

Un $(P, \xi)$ \struck{est dit} (relatif à $G$ admissible) est dit
\emph{minimal} si on ne peut trouver de sous-groupe admissible $G'$ de $G$, et
$(P', \xi')$ relatif à $G'$, t.q.

\[ (P, \xi) \simeq (P', \xi') \times_{G} G'. \]

\note{les deux groupes en indice et en facteur du produit sont lus comme ici ;
l'extension du groupe structural, du $G'$-fibré $(P',\xi')$ au $G$-fibré
$(P,\xi)$, demanderait $\times_{G'} G$ --- rapprocher de $(P,\xi) \times_{G}
(G' \times G')$ des pages 138 et 140, où l'indice est bien le groupe du fibré}

\page{138}

\struck{\ill{}} Soit $\mathfrak{Z}^{*}(a ; G) = $ ens.\ des classes d'isom.\ de
$G$-fibrés principaux \struck{\ill{}} pointés au-dessus de $a$. C'est un
\struck{\ill{}} covariant en le $S$-groupe admissible $G$.
\note{le $\mathfrak{Z}$ de ces feuillets est le $Z(S, a ; G)$ de la lettre des
pages 124 à 129}

\textbf{Proposition 4.} Si $\alpha \in \mathfrak{Z}(a, G)$ est \emph{minimal},
alors $u \mapsto u_{*}(\alpha)$ de $\operatorname{Hom}^{\mathrm{gr}}_{S}(G,
G')$ dans \struck{$\mathfrak{Z}$} $\mathfrak{Z}(a ; G')$ est \emph{injective}.

Soient en effet $u, v : G \to G'$ t.q.\ $u\alpha = v\alpha$. Soit $(P, \xi)$ de
classe $\alpha$, $(P', \xi')$ de classe $u\alpha = v\alpha$, $(u,v) : G \to G' \times G'$ (pris sur $S$)\note{l'indice du produit
$G' \times G'$ porte une rature}, d'où $(P, \xi) \times_{G} (G'
\times G')$, soit $(Q, \eta)$ ; par hypothèse, il y a un \ill{} compatible avec
les prolongations, et avec \add{la diagonale} $G' \to G' \times G'$ : $(P',
\xi') \to (P, \xi) \times_{G} (G' \times G')$, d'où un diagramme

\begin{tikzcd}
(P, \xi) \arrow[r] & (Q, \eta) \\
& (P', \xi') \arrow[u]
\end{tikzcd}

\begin{tikzcd}
G \arrow[r, "{(u,v)}"] & G' \times G' \\
& G' \arrow[u]
\end{tikzcd}

\page{139}

Soit $(R, \lambda)$ le produit fibré de $(P, \xi)$ et $(P', \xi')$ sur $(Q,
\eta)$,\note{la page s'arrête sur cette ligne ; l'argument est repris depuis le
début à la page 140}

\page{140}

$S$ préschéma connexe réduit

\struck{$\mathcal{C}$ une catégorie auxiliaire} \struck{$\mathcal{G}$ une
sous-catégorie pleine de $\mathrm{Gr}(\mathcal{C})$}

$\mathcal{C}$ une catégorie de $S$-groupes \emph{finis} et \emph{plats}.
$\mathcal{C}$ stable par \emph{produits}, \struck{et} par \emph{noyaux}
\struck{\ill{}} de couples de morphismes, par \emph{images}. Les $G \in
\mathcal{C}$ sont \emph{admissibles}.

$S$ \ill{} pointé par $a$.

Si $G \in \mathcal{C}$, on pose $\mathfrak{Z}(G) = \mathfrak{Z}(S, a ; G) = $
classes de fibrés principaux \ill{} $a$-pointés de groupe $G$.

Éléments \emph{minimaux} de $\mathfrak{Z}(G)$.

\textbf{Proposition 4.} Soit $\alpha \in \mathfrak{Z}(G)$ minimal. Alors $u
\mapsto u_{*}(\alpha)$ de $\operatorname{Hom}^{\mathrm{gr}}_{S}(G, G')$ dans
$\mathfrak{Z}(G')$ est \emph{injectif}. [$G$, $G' \in \mathcal{C}$.]

$u, v \in \operatorname{Hom}^{\mathrm{gr}}_{S}(G, G')$ tels que
\struck{\ill{}} $u_{*}(\alpha) = v_{*}(\alpha)$. $(P, \xi)$ de classe $\alpha$,
$(P', \xi')$ de classe $u_{*}(\alpha) = v_{*}(\alpha)$, $(u,v) : G \to G'
\times G'$, et

\[ (Q, \eta) = (P, \xi) \times_{G} (G' \times G') = [(P, \xi) \times_{u} G']
\times [(P, \xi) \times_{v} G'] \]

Par hypothèse, il y a un hom.\ $(P', \xi') \to (Q, \eta)$ (d'ailleurs unique)
compatible avec $G' \to G' \times G'$ (diagonale).

\begin{tikzcd}
G \arrow[r, "{(u,v)}"] & G' \times G' \\
H \arrow[u] \arrow[r] & G' \arrow[u]
\end{tikzcd}

\begin{tikzcd}
(P, \xi) \arrow[r] & (Q, \eta) \\
(R, \lambda) \arrow[u] \arrow[r] & (P', \xi') \arrow[u]
\end{tikzcd}

\note{dans le diagramme de gauche, chacune des lettres $G$, $G'$, $H$ porte un
trait horizontal que la page n'explique pas ; il n'est pas reproduit ici}

Soit $(R, \lambda)$ le produit fibré de $(P, \xi)$ et $(P', \xi')$ sur $(Q,
\eta)$. \struck{\ill{}} Soit $H$ le noyau de $G \rightrightarrows G'$ ; \ill{}
$(R, \lambda)$ est un fibré principal \emph{de groupe} $H$ \ill{}
\uncertain{et la minimalité de $(P, \xi)$} \ill{} donne $H = G$, d'où $u = v$.
En effet, il suffit de le montrer \struck{localement} en faisant l'extension de
base $(S_{1}, a_{1}) \to (S, a)$, ce qui amène à remplacer \ill{}

\marginal{$S \leftarrow S_{1}$, $g \mapsto g$ ; $u(g) = g'$, $v(g) = g_{1} g'$, $u(g) = g_{1}'^{-1} v(g)$}

\end{document}
